\documentclass{article}
\usepackage[margin=1in]{geometry}
\usepackage{mathtools, amsfonts, amssymb, amsthm, graphicx, listings, xcolor}


\definecolor{codegreen}{rgb}{0,0.6,0}
\definecolor{codegray}{rgb}{0.5,0.5,0.5}
\definecolor{codepurple}{rgb}{0.58,0,0.82}
\definecolor{backcolour}{rgb}{0.95,0.95,0.92}

\lstdefinestyle{mystyle}{
    backgroundcolor=\color{backcolour},   
    commentstyle=\color{codegreen},
    keywordstyle=\color{magenta},
    numberstyle=\tiny\color{codegray},
    stringstyle=\color{codepurple},
    basicstyle=\ttfamily\footnotesize,
    breakatwhitespace=false,         
    breaklines=true,                 
    captionpos=b,                    
    keepspaces=true,                 
    numbers=left,                    
    numbersep=5pt,                  
    showspaces=false,                
    showstringspaces=false,
    showtabs=false,                  
    tabsize=2
}

\lstset{style=mystyle}

\title{HW02}
\author{Sam Ly}

\begin{document}
\maketitle

Total: 27 pts (I just lost count lol)

\section*{[3 pts] Required Exercise 1.}

Done.

\begin{figure}[htbp]
    \centering
    \includegraphics[width=0.8\textwidth]{figures/prompt.png}
\end{figure}

\section*{[3 pts] Required Exercise 2.}

\begin{figure}[h]
    \centering
    \includegraphics[width=0.5\textwidth]{figures/hank.jpeg}
    \caption{Do I look like I know what a JPEG is?}
\end{figure}

Try the following settings for [width=...], and discuss what they do: 

\begin{enumerate}
    \item {
        \begin{verbatim}
            [width=2cm]
        \end{verbatim}

        This makes the figure very small, like a thumbnail. The
        sizing is abosulte w.r.t the literal page size.
    }
    \item {
        \begin{verbatim}
            [width=5cm]
        \end{verbatim}
        
        A bit bigger than the previous, but is still absolute sizing.
    }
    \item {
        \begin{verbatim}
            [width=0.5\textwidth]
        \end{verbatim}

        This size fills a good amount of the page without being overly 
        intrusive. I usually like using this.
    }

    \item {
        \begin{verbatim}
            [width=\linewidth]
        \end{verbatim}

        Seems to do the same as "textwidth". Internet says it has to do with 
        pages with multiple columns.
    }
\end{enumerate}

\section*{[4 pts] Required Exercise 3.}

\begin{enumerate}
    \item {
        The article took me 24 minutes to finish. I really like the raw, unfiltered
        voice of Jonathan Malesic. It seems he has been trying to articulate 
        these thoughts for a while. The part that I agree with most is his main 
        idea: that AI isn't the revolutionary tech that people seem to want to 
        believe. Sure it's impressive, and I also believe it is innovative, but 
        that's about it. It'll change specific aspects of certain jobs,
        software engineering per se, but it is clear to see that it can't replicate
        the core of "cultural output": the humanness of ideas.

        At it's core, AI is a fancied-up fuzzy search. Even recent breakthroughs
        that claim to push this envelope like the Reasoning Model have been shown 
        to actually be memorizing old patterns. In other words, a fuzzy search. 
        The "generative" in Gen AI is really just searching through old content 
        in hopes it'll find something that matches your prompt. 

        Now one thing that I disagree with is the notion that AI is just a gimmick. 
        It does provide value in a subset of problems that require broad, high-level
        knowledge and understanding. Although standard AI chat bots falls short 
        in tasks requiring deep, domain-specific knowledge, it does great in giving
        you a "place to start" in your search. I've experienced this first hand 
        in a lot of software engineering scenarios (I recognize the potential for
        bias here), where I almost always use AI to build project templates, 
        evaluate pros and cons of libraries/frameworks, and generate boilerplate. 

        Questions that used to require years of experience to answer effectively
        can now be answered "good-enough-ly" by a chat bot. This is where the true 
        value of AI in engineering lies, acting has a heuristic for high level decisions.
        Again, AI design decisions can and will be wrong a percentage of the time,
        but it is folly to assume the same errors never occur with human engineers.
    
    }

    \item {
        My favorite quote from the article is "AI progress in cultural production 
        already seems to have slowed because the models have run out of 
        human-generated writing to “learn” from and increasingly feed on AI-produced 
        content, gulping down a vile soup of their own ever-concentrating ordure." 

        This quote really resonates with me because I hate seeing AI slop online. 
        It seems everywhere I look I see it. I can't avoid it, and so the best I 
        can do is ignore it. From YouTube shorts featuring and AI voice over 
        narrating a nonsensical AI generated story, to emoji-filled LinkedIn 
        "guides", AI generated content is filling the internet. 

        If I wanted a AI generated response, I would prompt it myself. It is, in 
        my opinion, at least slightly disrespectful to post AI content like 
        this because it is essentially saying that you don't think the audience 
        is smart enough to know.
    }
\end{enumerate}

\section*{[9 pts] Choice Exercise 4.}

\begin{enumerate}
    \item {
        Prove that given 5¢ coins and 6¢ coins, you can make change for any 
        amount of money that is 20¢ or more.

        \begin{proof}
            We begin by defining "being able to make change from 5¢ and 6¢ coins"
            as:
            \[n = 5a + 6b\]
            
            Now, we see that we can make change of \(n\)¢ if we can write \(n\)
            as \(5a + 6b\). 

            For the base cases \(20 \le n \le 24\), we have:
            \begin{itemize}
                \item \(20 = 5(4) + 6(0)\).
                \item \(21 = 5(3) + 6(1)\).
                \item \(22 = 5(2) + 6(2)\).
                \item \(23 = 5(1) + 6(3)\).
                \item \(24 = 5(0) + 6(4)\).
            \end{itemize}

            Now we make the inductive hypothesis that we \textit{can} make 
            change of any value \(n = k\) such that \(k \ge 20\). It's clear 
            that cases \(20 \le k \le 24\) is equivalent to the above base cases, 
            and so the relation holds.

            We can also make change for \(k \ge 25\), because 
            \(k + 1 = k - 4 + 5\). We can see this by taking the specific case
            \(k = 25\):
            \[25 = 24 + 1 = 20 + 5 = 5(4) + 5 = 5(5).\]

            Similar logic applies to all \(k \ge 25\).

            Therefore, we can make change for any amount of money that is 20¢ or 
            more using only 5¢ and 6¢ coins.
        \end{proof}
    }

    \item {
        Prove that \(9^n + 5^n - 2\) is divisible by 4 for all integers \(n \ge 1\). 

        \begin{proof}
            An integer \(m\) is divisible by 4 if it can be written in the form 
            \(m = 4c\) for some integer \(c\).

            For \(n=1\), we have \(9^1 + 5^1 - 2 = 12 = 4(3)\).

            For \(n=2\), we have \(9^2 + 5^2 - 2 = 104 = 4(26)\).

            We can make the inductive hypothesis that \(9^n + 5^n - 2 = 4m\)
            holds true for \(1 \le n \le k\).

            Thus,
            \[9^k = 4m - 5^k + 2\]
            and
            \[5^k = 4m - 9^k + 2.\]

            Then, we see if the relation holds for \(n = k+1\) via the following 
            algebraic manipulation: 
            \[9^{k+1} + 5^{k+1} - 2\]
            \[9(9^k) + 5(5^k) - 2\]
            \[9(4m - 5^k + 2) + 5(4m - 9^k + 2) - 2\]
            \[36m - 9(5^k) + 18 + 20m - 5(9^k) + 10 - 2\]
            \[56m - 9(5^k) + 5(9^k) + 26\]
            \[56m - 45(\frac{5^k}{5} + \frac{9^k}{9}) + 26\]
            \[56m - 45(5^{k-1} + 9^{k-1} - 2) - 2(45) + 26\]
            \[56m - 45(5^{k-1} + 9^{k-1} - 2) - 64.\]
            
            Now, we can use our inductive hypothesis to say 
            \(5^{k-1} + 9^{k-1} - 2 = 4m_1\).

            Thus we have 
            \[9^{k+1} + 5^{k+1} - 2 = 56m - 45(4m_1) - 64 = 4(14m - 45m_1 - 16).\]

            Therefore, \(9^n + 5^n - 2\) is divisible by 4 for all \(n \ge 1\).
        \end{proof}
    }

    \item {
        Let \(F_n\) be the \(n\)-th Fibonacci number, where \(F_0 = F_1 = 1\)
        and \(F_n = F_{n-1} + F_{n-2}\). Prove that \(F_n \le 1.9^n\) for all
        \(n \ge 1\).

        \begin{proof}
            We begin by verifying the relation for small \(n\) to create our 
            base cases:
            \[n = 1, F_1 = 1 \le 1.9^1 = 1.9\]
            \[n = 2, F_2 = 2 \le 1.9^2 = 3.61.\]

            We then form our inductive hypothesis by assumimg \(F_n \le 1.9^n\) for \(1 \le n \le k\).
            \[F_{k+1} \le 1.9^{k+1}\]
            \[F_k + F_{k-1} \le 1.9^{k+1}.\]

            Using our inductive hypothesis, we see that \(F_k \le 1.9^k\)
            and \(F_{k-1} \le 1.9^{k-1}\).

            So, \(F_k + F_{k-1} \le 1.9^k + 1.9^{k-1}\).

            By refactoring \(1.9^k + 1.9^{k-1}\), we get:
            \[1.9^k + 1.9^{k-1} = 1.9(1.9^{k-1}) + 1.9^{k-1} = 2.9(1.9^{k-1}).\]

            Also, \(1.9^{k+1}\) can be rewritten as \(1.9^2 (1.9^{k-1}) = 3.61(1.9^{k-1})\).

            Finally, we see
            \[F_{k+1} = F_k + F_{k-1} \le 2.9(1.9^{k-1}) \le 3.61(1.9^{k-1}) = 1.9^{k+1}\]
            \[F_{k+1} \le 1.9^{k+1}.\]

            Therefore, \(F_n \le 1.9^n\) for all \(n \ge 1\).
        \end{proof}
    }
\end{enumerate}

\section*{[3 pts] Choice Exercise 5.}

\begin{enumerate}
    \item[3.] {
        Prove or disprove that if \(n\) is a positive integer, then the concatenation 
        \(nnn\) is divisible by \(37\). 

        \begin{proof}
            We first define the concatenation function \(C\) as: 
            \[C(n) = \begin{cases}
                111 \times n & 1 \le n \le 9 \\
                10101 \times n & 10 \le n \le 99 \\
                1001001 \times n & 100 \le n \le 999 \\
                ...

            \end{cases}\]

            For the cases \(1 \le n \le 9\) and \(10 \le n \le 99\), our proposition 
            holds true since \(37 \mid 111\) and \(37 \mid 10101\). Thus, \(37 \mid C(n)\)
            for \(1 \le n \le 99\).

            However, \(37 \nmid 1001001\), so there may be a case \(100 \le n \le 999\) 
            where \(37 \nmid C(n)\). 

            Since \(999 - 100 > 37\), by the Pigeon Hole Principle (maybe, I'm 
            too sure if I am using this Principle correctly), there must be a 
            value \(100 \le n \le 999\) where \(37 \nmid n\). 

            By extension, because 

            Thus, we have disproven the claim that the concatenation of any positive
            integer \(n\) is divisible by \(37\). 

            For example, the concatenation of 186 is 186,186,186, which is not 
            divisible by 37.
        \end{proof}
    }
\end{enumerate}

\section*{[5 pts] Choice Exercise 8.}

The error in this proof comes at the inductive step, where we test if the claim 
holds true for \(n = k+1\). 

Remember, our inductive hypothesis is meant to say, "If all the cases \(n = 1, 2, ..., k\) 
are true, then the case \(n = k+1\) must also be true." We must verify this statement
via algebra, or what ever else mathematical tools. 

This proof fails in this regard as it does not verify that the truth of cases 
\(n = 1, 2, ..., k\), \textbf{implies} the truth of case \(n = k+1\). It instead 
\textit{assumes} that case \(n = k+1\) is true. When doing a proof by induction,
we can only assume the truth of \(n = 1, 2, ..., k\), then \textit{verify} that 
this is the "sufficicient condition" for case \(n = k+1\).

In other words, the inductive hypothesis can only extend up to case \(n = k\).
Case \(n = k+1\) must be verified. We can not simply assume that the \(k+1\)-th 
horse is the same color as the \(k\)-th horse. Thus this proof is invalid.

\end{document}